\documentclass[10pt]{article}
\usepackage[utf8]{inputenc}
\usepackage[T1]{fontenc}
\usepackage{geometry}
\usepackage{booktabs}
\usepackage{multirow}
\usepackage{amsmath,amssymb}
\usepackage{graphicx}
\usepackage{hyperref}
\geometry{margin=1in}

\title{Mechanisms of Social Status Regulation as Free--Energy Minimising Loops}
\author{Compiled for AI Safety Collaborators}
\date{August 2025}

\begin{document}
\maketitle

\begin{abstract}
Human status dynamics appear to be governed by a portfolio of partially independent optimisation loops, each minimising a distinct component of variational free energy. We propose a compact taxonomy---from reptilian dominance reflexes to language--mediated reputation markets---and specify, for each loop, (i) the free--energy term measured, (ii) the neuroanatomical registers and actuators, and (iii) a plausible micro--mechanism that yields a prediction--error (PE) signal with neuromodulatory precision--gating. We articulate an internal/external symmetry (interoceptive threat vs. exteroceptive perimeter defence; somatic repair vs. niche construction; etc.), derive perturbation predictions, and outline how these principles transfer to opaque--agent forensics and brain--like AI benchmarks.
\end{abstract}

\section{Background}
Under the free--energy principle (FEP) \cite{Friston2010}, adaptive systems minimise expected surprise by building predictive models and selecting actions that make sensory data conform to those models. Status regulation can be cast as a set of social free--energy ledgers: uncertainty about threat, coalition membership, norm adherence, prestige sources, and reputational standing. Evolutionary continuity suggests layered solutions: brainstem circuits for dyadic dominance, limbic coalition bookkeeping, and neocortical narrative broadcast \cite{Sapolsky2004,HenrichGilWhite2001,Zink2008}.

\section{Loops, ledgers, and evolutionary stratification}
Table~\ref{tab:ledger} orders status loops from older to newer and states, for each, the free--energy target $F$, registers, actuators, and neuromodulators that set PE precision.

\begin{table}[h]
\small
\centering
\begin{tabular}{@{}r p{2.6cm} p{3.6cm} p{3.6cm} p{2.1cm} p{1.9cm}@{}}
\toprule
Age & Loop & Free--energy term $F$ & Registers (priors / evidence) & Actuators & Precision gate \\
\midrule
1 & Dominance--submission & $F_{\mathrm{threat}}=\Pr(\mathrm{injury}\mid \pi)$ & BLA prior \,$\leftrightarrow$\, VMHvl, MeA evidence & PAG columns (fight/flight), spinal patterns & 5--HT, T \cite{Lin2011,Raleigh1991} \\
2 & Display--assessment & $F_{\mathrm{signal}}=\lvert\hat s_{\mathrm{self}}-s_{\mathrm{peer}}\rvert$ & Ventral striatum template \,$\leftrightarrow$\, auditory/visual cortex & Pattern generators (laryngeal, facial) & DA \cite{Schultz1997} \\
3 & Coalition entropy & $H(C_{ij})$ & TPJ/mPFC graph prior \,$\leftrightarrow$\, STS cues & Grooming/approach, alliance choice & OXT \cite{Scholz2016,Kosfeld2005} \\
4 & Prestige--biased learning & $F_{\mathrm{ep}}^{\mathrm{social}}=\mathbb E\,\lvert o-\hat o_{\mathrm{model}}\rvert$ & IFG+pSTS mirror prior \,$\leftrightarrow$\, PPC/cerebellar mismatch & Imitation weighting, teaching & DA \cite{Rizzolatti2004,Chudek2011} \\
5 & Norm enforcement & $F_{\mathrm{norm}}=\Pr(\mathrm{rule\ breach})$ & vmPFC/hippocampal rule prior \,$\leftrightarrow$\, ACC/insula error & dlPFC$\to$BG punishment policy & $\mu$--opioid, NA \cite{Fehr2002,Knoch2006,Eisenberger2003} \\
6 & Reputation broadcast & $F_{\mathrm{rep}}=\lvert \hat r_i-r_i\rvert$ & DMN (PCC, temp. pole) prior \,$\leftrightarrow$\, STG language input & Narrative update, signalling & DA/\,OXT \cite{Zink2008,HenrichGilWhite2001} \\
\bottomrule
\end{tabular}
\caption{Status loops from evolutionarily older (1) to newer (6). Abbreviations: BLA basolateral amygdala; VMHvl ventromedial hypothalamus (ventrolateral); MeA medial amygdala; TPJ temporoparietal junction; STS superior temporal sulcus; IFG inferior frontal gyrus; PPC posterior parietal cortex; dlPFC dorsolateral prefrontal cortex; BG basal ganglia; DMN default mode network; NA noradrenaline; OXT oxytocin; DA dopamine; 5--HT serotonin; T testosterone.}
\label{tab:ledger}
\end{table}

\section{How circuits \emph{measure} the ledgers}
Across loops, a common micro--architecture appears: deep pyramidal cells carry priors; superficial cells compute residuals; neuromodulators scale precision (inverse variance). Concretely:
\begin{itemize}
  \item \textbf{Dominance threat:} BLA sends expected threat to VMHvl; MeA relays opponent cues. Subtractive interneurons in VMHvl produce a PE that recruits PAG; raphe 5--HT sets risk tolerance, testosterone gates plasticity \cite{Lin2011,Sapolsky2004}.
  \item \textbf{Coalition entropy:} TPJ/mPFC encode a sparse triadic graph prior. pSTS supplies live affiliative cues. Error neurons in TPJ fire when link likelihood deviates, with OXT increasing gain (precision) on in--group edges \cite{Scholz2016,Kosfeld2005}.
  \item \textbf{Prestige:} Mirror system (IFG/pSTS) hosts a predictive motor program of the model; cerebellum/PPC compute kinematic mismatch; VTA DA encodes epistemic PE that reweights imitation \cite{Rizzolatti2004,Schultz1997}.
  \item \textbf{Norms:} vmPFC stores rule templates; ACC compares allowed action vs. observed; an aversive PE drives dlPFC$\to$BG sanction. Relief upon resolution corresponds to a drop in $F_{\mathrm{norm}}$ and engages $\mu$--opioid \cite{Knoch2006,Eisenberger2003}.
  \item \textbf{Reputation:} DMN maintains a semantic centrality prior; language cortex updates evidence; ventral striatum emits a value PE proportional to change in expected future cooperation \cite{Zink2008}.
\end{itemize}

\section{Internal $\leftrightarrow$ external symmetry}
Each interoceptive ledger has an exteroceptive twin: internal threat (amygdala/PAG) vs. external perimeter defence (SC--guided vigilance); internal somatic repair vs. external niche construction; internal norm adherence (guilt) vs. external sanction (punishment). This symmetry follows from FEP at two scales: body maintenance and niche shaping.

\section{Perturbation predictions}
\subsection*{TMS to dlPFC (norm loop)}
Continuous theta--burst TMS over right dlPFC should attenuate third--party punishment without reducing ACC mismatch (N200/P300). Behaviourally, acceptance of unfair splits rises; sanction spending falls \cite{Knoch2006}. Relief signals (opioid) diminish until excitability recovers.

\subsection*{Serotonergic shift (dominance loop)}
An acute SSRI dose (increased 5--HT) raises the submit threshold and reduces aggressive approach under status threat; the PE amplitude in VMHvl drops for identical opponent cues \cite{Raleigh1991}.

\section{Consequences for agency and consciousness}
Status loops feed into valuation--efference circuits: changes in $F_{\mathrm{threat}}$ and $H(C_{ij})$ directly alter discounting, exploration, and social policy selection \cite{Daw2006}. Global broadcast and metacognition render these computations available to report, shaping conscious feelings of pride, shame, outrage. Miscalibrated precision (e.g., high OXT precision on in--group edges) can lower social free energy locally while increasing global risk (out--group hostility), a classical alignment pitfall for artificial agents.

\section{Opaque--agent forensics: low--sample discovery}
Borrowing from lesion logic, we recommend four probes per discovered module: (i) topology screen (hub vs. chain), (ii) one--bit poke (rebound gain $g>1$ indicates homeostasis), (iii) 30~s behavioural battery (rank cue, norm breach, prestige honeypot), and (iv) power snapshot. A triple--hit (structure$+$PE$+$behaviour) labels a status loop without long recordings.

\section{Outlook}
The status stack recapitulates vertebrate evolution while exploiting neocortical broadcast. Incorporating these ledgers into AI benchmarks should surface prestige--seeking, coalition maintenance, and norm enforcement as distinct optimisation targets, enabling targeted ``cooling'' or constraint of each loop.

\paragraph{Limitations.} Citations are indicative; multiple parallel routes can implement the same PE computation. Quantitative priors (precision) are context--dependent and ethically sensitive to manipulate in humans.

\begin{thebibliography}{99}
\bibitem{Friston2010} K. Friston (2010). The free-energy principle: a unified brain theory? \emph{Nat. Rev. Neurosci.} 11:127--138.
\bibitem{Sapolsky2004} R. Sapolsky (2004). Social status and health in humans and other animals. \emph{Annu. Rev. Anthropol.} 33:393--418.
\bibitem{HenrichGilWhite2001} J. Henrich \& F. Gil-White (2001). The evolution of prestige. \emph{Evol. Hum. Behav.} 22:165--196.
\bibitem{Zink2008} C. F. Zink et al. (2008). Know your place: neural processing of social hierarchy. \emph{Neuron} 58:273--283.
\bibitem{Lin2011} D. Lin et al. (2011). Functional identification of an aggression locus in the mouse hypothalamus. \emph{Nature} 470:221--226.
\bibitem{Raleigh1991} M. J. Raleigh et al. (1991). Serotonergic mechanisms in vervet dominance. \emph{Brain Res.} 559:181--190.
\bibitem{Schultz1997} W. Schultz et al. (1997). A neural substrate of prediction and reward. \emph{Science} 275:1593--1599.
\bibitem{Scholz2016} J. Scholz et al. (2016). rTPJ in attention and social cognition: combined fMRI--TMS. \emph{Cereb. Cortex} 26:4282--4291.
\bibitem{Kosfeld2005} M. Kosfeld et al. (2005). Oxytocin increases trust. \emph{Nature} 435:673--676.
\bibitem{Rizzolatti2004} G. Rizzolatti \& L. Craighero (2004). The mirror-neuron system. \emph{Annu. Rev. Neurosci.} 27:169--192.
\bibitem{Chudek2011} M. Chudek \& J. Henrich (2011). Culture--gene coevolution, prestige and conformity. \emph{Phil. Trans. R. Soc. B} 366:1149--1158.
\bibitem{Fehr2002} E. Fehr \& S. G\"achter (2002). Altruistic punishment. \emph{Nature} 415:137--140.
\bibitem{Knoch2006} D. Knoch et al. (2006). Disruption of right dlPFC decreases norm compliance. \emph{Science} 314:829--832.
\bibitem{Eisenberger2003} N. I. Eisenberger et al. (2003). Does rejection hurt? \emph{Science} 302:290--292.
\bibitem{Daw2006} N. D. Daw et al. (2006). Cortical substrates for model-based vs. model-free learning. \emph{Nat. Neurosci.} 8:1704--1711.
\end{thebibliography}

\end{document}
